\documentclass[11pt]{amsart}
\usepackage{amsmath,amssymb,amsthm, mathtools}
\usepackage[margin=1in]{geometry}
\usepackage{enumitem}

\pagestyle{plain}

\begin{document}

\begin{center}
    {\large Math 4101 Practice Midterm}\\[2pt]
    {\large Extra Practice Set}
\end{center}

\noindent Name: \underline{\hspace{2in}}

\medskip
\noindent Time: 50 minutes \hfill Total: 100 points

\medskip
\noindent \textbf{Directions.} Write complete proofs or arguments. Clearly state any definitions you use.
Unless otherwise stated, all numbers are real. Throughout this exam, $\mathbb{N} = \{1, 2, 3, \dots\}$.

\bigskip

\noindent\textbf{Problem 1 (30 points).}

\noindent Let
\[
S = \{x \in \mathbb{R} : x^3 \in \mathbb{Q}\}.
\]

\begin{enumerate}[label=(\alph*)]
    \item (15 points) Prove that $S$ is countably infinite.
    \item (15 points) Prove that $\mathbb{R} \setminus S$ is uncountable.
\end{enumerate}

\bigskip

\noindent\textbf{Problem 2 (30 points).}

\noindent Let $A \subseteq \mathbb{R}$ be nonempty and bounded below, and write $m = \inf A$.

\begin{enumerate}[label=(\alph*)]
    \item (15 points) Prove that for every $\varepsilon > 0$, there exists $a \in A$ such that
    \[
    m \le a < m + \varepsilon.
    \]
    \item (15 points) Define
    \[
    T = \left\{ a + \frac{1}{n} : a \in A,\ n \in \mathbb{N} \right\}.
    \]
    Prove that $\inf T = m$. You may use part (a), even if you have not proved it.
\end{enumerate}

\bigskip

\noindent\textbf{Problem 3 (40 points).}

\noindent Use the standard metric $d(x,y) = |x-y|$ on $\mathbb{R}$.

\begin{enumerate}[label=(\alph*)]
    \item (20 points) Let $U, V \subseteq \mathbb{R}$ be open. Prove directly from the definition of openness that $U \cup V$ is open.
    \item (20 points) Suppose $U_n \subseteq \mathbb{R}$ is open for every $n \in \mathbb{N}$. Must
    \[
    \bigcup_{n=1}^{\infty} U_n
    \]
    be open? Give a proof or an explicit counterexample. If you give a counterexample, determine the union and justify your answer.
\end{enumerate}

\bigskip

\noindent\textbf{Problem 4 (Bonus, 10 points).}

\noindent Suppose $F_n \subseteq \mathbb{R}$ is closed for every $n \in \mathbb{N}$. Must
\[
\bigcup_{n=1}^{\infty} F_n
\]
be closed? Give a proof or an explicit counterexample. (Hint: think about how this problem relates to Problem 3, part (b), via complements and De Morgan's laws.)

\end{document}